% Part: first-order-logic
% Chapter: beyond
% Section: higher-order-logic

\documentclass[../../../include/open-logic-section]{subfiles}

\begin{document}

% SEGMENT OLP-0178-S01

\olfileid{fol}{byd}{hol}

\olsection{உயர்வரிசைத் தருக்கம்}

முதல்வரிசைத் தருக்கத்திலிருந்து இரண்டாம்வரிசைத் தருக்கத்திற்குச் சென்றது,
முறைசார் மொழிக்குள்ளேயே முதல்வரிசை !!{domain} இலுள்ள பொருட்களின்
கணங்களைப் பற்றிப் பேச நமக்கு வழிவகுத்தது. ஏன் அத்துடன் நிறுத்த வேண்டும்?
எடுத்துக்காட்டாக, மூன்றாம்வரிசைத் தருக்கம் பொருட்களின் கணங்களின்
கணங்களையோ, பொருட்களையும் பொருட்களின் கணங்களையும் ஒருங்கே கொண்ட
கணங்களையோகூட கையாள நமக்கு வழிவகுக்க வேண்டும். நான்காம்வரிசைத் தருக்கம்
அவ்வகைப் பொருட்களின் கணங்களைப் பற்றிப் பேச அனுமதிக்கும். நீங்கள்
ஊகித்திருக்கக்கூடியபடி, இக்கருத்தை எத்தனை முறையும் தொடரலாம்.

நடைமுறையில், உயர்வரிசைத் தருக்கம் தொடர்புகளுக்குப் பதிலாகச் சார்புகளைக்
கொண்டு அடிக்கடி !!{formula}க்கப்படுகிறது. (இயல்பான ஒன்றாக்கங்களைப்
பொறுத்தவரை, இவ்வேறுபாடு சாரமற்றது.) சில அடிப்படை ``வகைகள்'' $A$, $B$,
$C$,~\dots (இனி இவற்றை ``வகைப்பாடுகள்'' என்று அழைப்போம்) கொடுக்கப்பட்டால்,
பின்வருமாறு விதித்து புதியவற்றை உருவாக்கலாம்:
\begin{quote}
$\sigma$ மற்றும் $\tau$ இறுதிக்குட்பட்ட வகைப்பாடுகள் எனில்,
$\sigma \to \tau$ உம் அவ்வாறே.
\end{quote}
வகைப்பாடுகளை, நமது !!{domain} இல் வேண்டிய பொருட்களை வகைப்படுத்தும்
தொடரியல் ``சீட்டுகள்'' எனக் கருதுக; வகைப்பாடு~$\sigma$ இன்
பொருட்களை வகைப்பாடு~$\tau$ இன் பொருட்களுக்கு எடுத்துச்செல்லும்
சார்புகளான பொருட்களை $\sigma \to \tau$ விவரிக்கிறது. எடுத்துக்காட்டாக,
``மெய்'', ``பொய்'' என்ற உண்மை மதிப்புகளின் வகைப்பாடு~$\Omega$ ஒன்றையும்,
இயல் எண்களின் வகைப்பாடு~$\Nat$ ஒன்றையும் நாம் கொள்ளலாம். அப்போது
$\Nat \to \Omega$ வகைப்பாட்டின் பொருட்களை ஓரிடத் தொடர்புகளாகவோ
$\Nat$ இன் உட்கணங்களாகவோ கருதலாம்; $\Nat \to \Nat$ வகைப்பாட்டின்
பொருட்கள் இயல் எண்களிலிருந்து இயல் எண்களுக்கான சார்புகள்; மேலும்
$(\Nat \to \Nat) \to \Nat$ வகைப்பாட்டின் பொருட்கள் ``சார்பிகள்'',
அதாவது சார்புகளை எண்களுக்கு எடுத்துச்செல்லும் உயர்வகைப்பாட்டுச் சார்புகள்.

இரண்டாம்வரிசைத் தருக்கத்தைப் போலவே, நாம் கருத விரும்பும் ஒவ்வொரு
வகைப்பாட்டுப் பொருளுக்கும் ஒரு வகை உள்ள பல-வகைத் தருக்கமாக
உயர்வரிசைத் தருக்கத்தைக் கருதலாம். ஆனால் உயர்வகைப்பாட்டுத் தருக்கத்தின்
தொடரியலை அடிப்படையிலிருந்தே வரையறுப்பது பொதுவாகத் தெளிவானது.
எடுத்துக்காட்டாக, இறுதிக்குட்பட்ட வகைப்பாடுகளின் கணத்தைப் பின்வருமாறு
தொகுத்தறிவாக வரையறுக்கலாம்:
\begin{enumerate}
\item $\Nat$ ஓர் இறுதிக்குட்பட்ட வகைப்பாடு.
\item $\sigma$ மற்றும் $\tau$ இறுதிக்குட்பட்ட வகைப்பாடுகள் எனில்,
  $\sigma \to \tau$ உம் அவ்வாறே.
\item $\sigma$ மற்றும் $\tau$ இறுதிக்குட்பட்ட வகைப்பாடுகள் எனில்,
  $\sigma \times \tau$ உம் அவ்வாறே.
\end{enumerate}
உள்ளுணர்வாக, $\Nat$ இயல் எண்களின் வகைப்பாட்டைக் குறிக்கிறது;
$\sigma \to \tau$ ஆனது $\sigma$ இலிருந்து $\tau$ க்கான சார்புகளின்
வகைப்பாட்டையும், $\sigma \times \tau$ ஆனது $\sigma$ இலிருந்து ஒன்றும்
$\tau$ இலிருந்து ஒன்றுமாக இரு பொருட்களின் சோடிகளின் வகைப்பாட்டையும்
குறிக்கின்றன. பின்னர் சொற்களின் கணத்தைப் பின்வருமாறு தொகுத்தறிவாக
வரையறுக்கலாம்:
\begin{enumerate}
\item ஒவ்வொரு வகைப்பாடு~$\sigma$ க்கும், அவ்வகைப்பாட்டின்
  மாறிகள் $x$, $y$, $z$, \dots ஆகியவற்றின் கையிருப்பு உண்டு; அவை
  வகைப்பாடு $\sigma$ இற்குரியவை.
\item $\Obj 0$ ஆனது $\Nat$ வகைப்பாட்டின் ஒரு சொல்.
\item $\Obj S$ (அடுத்தெண்) ஆனது $\Nat \to \Nat$ வகைப்பாட்டின் ஒரு சொல்.
\item $s$ ஆனது $\sigma$ வகைப்பாட்டின் சொல்லாகவும், $t$ ஆனது
  $\Nat \to (\sigma \to \sigma)$ வகைப்பாட்டின் சொல்லாகவும் இருந்தால்,
  $\Obj{R}_{st}$ ஆனது $\Nat \to \sigma$ வகைப்பாட்டின் ஒரு சொல்.
\item $s$ ஆனது $\tau \to \sigma$ வகைப்பாட்டின் சொல்லாகவும், $t$ ஆனது
  வகைப்பாடு~$\tau$ இன் சொல்லாகவும் இருந்தால், $s(t)$ ஆனது
  $\sigma$ வகைப்பாட்டின் ஒரு சொல்.
\item $s$ ஆனது வகைப்பாடு~$\sigma$ இன் சொல்லாகவும், $x$ ஆனது
  வகைப்பாடு~$\tau$ இன் மாறியாகவும் இருந்தால், $\lambd[x][s]$ ஆனது
  $\tau \to \sigma$ வகைப்பாட்டின் ஒரு சொல்.
\item $s$ ஆனது வகைப்பாடு~$\sigma$ இன் சொல்லாகவும், $t$ ஆனது
  வகைப்பாடு~$\tau$ இன் சொல்லாகவும் இருந்தால், $\tuple{s, t}$ ஆனது
  $\sigma \times \tau$ வகைப்பாட்டின் ஒரு சொல்.
\item $s$ ஆனது வகைப்பாடு~$\sigma \times \tau$ இன் சொல்லாக இருந்தால்,
  $p_1(s)$ ஆனது வகைப்பாடு~$\sigma$ இன் சொல்லும், $p_2(s)$ ஆனது
  வகைப்பாடு~$\tau$ இன் சொல்லும் ஆகும்.
\end{enumerate}
உள்ளுணர்வாக, $\Obj{R}_{st}$ ஆனது பின்வருமாறு மீள்வரையறுக்கப்பட்ட
சார்பைக் குறிக்கிறது:
\begin{align*}
\Obj{R}_{st}(0) & = s \\
\Obj{R}_{st}(x+1) & = t(x, R_{st}(x)),
\end{align*}
$\tuple{s, t}$ ஆனது முதல் கூறு~$s$ ஆகவும் இரண்டாம் கூறு~$t$ ஆகவும்
உள்ள சோடியைக் குறிக்கிறது; $p_1(s)$ மற்றும்~$p_2(s)$ ஆகியவை~$s$ இன்
முதல், இரண்டாம் உறுப்புகளை (``வீழ்த்தல்களை'') குறிக்கின்றன. இறுதியாக,
$\lambd[x][s]$ ஆனது வகைப்பாடு~$\sigma$ இன் ஏதேனும்~$x$ க்கும்
\[
f(x) = s
\]
என்று வரையறுக்கப்படும் சார்பு~$f$ ஐக் குறிக்கிறது; ஆகவே உருப்படி (6)
சொற்களைக் கொண்டு சார்புகளை வரையறுக்க நமக்கு வழிவகுக்கும் உட்கொள்ளலின்
ஒரு வடிவத்தைத் தருகிறது. ஒரே வகைப்பாட்டுச் சொற்களுக்கிடையிலான
!!{identity} கூற்றுகள் $\eq[s][t]$, வழக்கமான கூற்றுக்கணிப்பு
இணைப்பிகள், உயர்வகைப்பாட்டு அளவையீடு ஆகியவற்றிலிருந்து
!!^{formula}s கட்டப்படுகின்றன. மேலே வரையறுக்கப்பட்ட சொற்களை ஆளும்
அடிப்படைச் சமன்பாடுகளையும், அளவையீட்டுக்குறிகள் மற்றும் !!{identity}
உடைய வழக்கமான தருக்க விதிகளையும் அமைப்பின் அடிகோள்களாகக் கொள்ளலாம்.

இறுதிக்குட்பட்ட வகைப்பாட்டு அமைப்புடன் உண்மை மதிப்புகளின்
வகைப்பாடு~$\Omega$ ஒன்றைச் சேர்த்தால், அதன் பயன்பாட்டை ஆளும்
அடிகோள்களையும் சேர்க்க வேண்டும். உண்மையில், திறமையாகச் செய்தால்,
சிக்கலான !!{formula}s ஐ முற்றிலும் நீக்கி, அவற்றை வகைப்பாடு~$\Omega$
இன் சொற்களால் மாற்றலாம்!{} அதற்கேற்ப நிரூபிப்பு அமைப்பையும் மாற்றலாம்.
இதன் விளைவு, 1930களில் அலோன்சோ சர்ச் முன்வைத்த
\emph{வகைப்பாடுகளின் எளிய கோட்பாடு} சாரமாகும்.

இரண்டாம்வரிசைத் தருக்கத்தைப் போலவே, பயன்படுத்த விரும்பத்தக்க
உயர்வகைப்பாட்டுப் பொருண்மையியலின் வெவ்வேறு வடிவங்கள் உண்டு. முழு
வடிவத்தில், $\sigma \to \tau$ வகைப்பாட்டின் மாறிகள்
வகைப்பாடு~$\sigma$ இன் பொருட்களிலிருந்து வகைப்பாடு~$\tau$ இன்
பொருட்களுக்கான \emph{எல்லா} சார்புகளின் கணத்தின் மீதும் வீச்சுறுகின்றன.
எதிர்பார்க்கக்கூடியபடி, இப்பொருண்மையியல் நிறைவான, செயற்படத்தக்க
!!{derivation} அமைப்பை அனுமதிக்க முடியாத அளவுக்கு வலுவானது. ஆனால்,
ஒவ்வொரு வகைப்பாடு $\tau$ க்குமான உறுப்புகளின் கணங்கள் $T_\tau$ உடனும்,
சார்புப் பயன்பாடு, வீழ்த்தல் முதலியவற்றுக்கான உரிய செயல்களுடனும் சேர்ந்த
!!a{structure} கொண்ட வலுக்குறைந்த பொருண்மையியலைக் கருதலாம். விவரங்களைச்
சரியாக அமைத்தால், மேலே விவரித்த வகை !!{derivation} அமைப்புகளுக்கான
நிறைவுத் தேற்றங்களைப் பெறலாம்.

கணிதத்தின் பெரும்பகுதியை இயல்பான முறையில் பதிக்கக்கூடிய சட்டகத்தை
உயர்வகைப்பாட்டுத் தருக்கம் வழங்குவதால் அது ஈர்ப்புடையது: $\Nat$ இல்
தொடங்கி, மெய்யெண்கள், தொடர்ச்சியான சார்புகள் முதலியவற்றை வரையறுக்கலாம்.
வகைப்பாடுகளுக்குத் தெளிவான ``கட்டுமான'' பொருள்கோள்கள் இருப்பதால்,
உள்ளுணர்வுவாதத் தருக்கச் சூழலிலும் இது குறிப்பாக ஈர்ப்புடையது.
உண்மையில், இக்கட்டுமானப் பொருள்கோள்களை மிகத் தெளிவாக விளக்கும்
உயர்வகைப்பாட்டுப் பொருண்மையியலின் கட்டுமான வடிவங்களை (செவ்வியல்
தருக்கத்திற்குப் பதிலாக உள்ளுணர்வுவாதத் தருக்கத்தைச் சார்ந்தவற்றை)
உருவாக்கலாம்; பலவிதங்களில் அவை செவ்வியல் இணையுருக்களைவிடச்
சுவையானவையாகும்.

\end{document}
